\section*{Chapitre \ref{ch:b3:submanifolds} --- Sous-variétés de
$\R^n$}

\begin{solution}{exo:b3:submanifolds:1}
(a) Chacun est $F^{-1}(0)$ pour une submersion~: $\norm x^2 -
1$ sur $\R^n\setminus\{0\}$ (gradient $2x \neq 0$)~:
dimension $n-1$~; $x^2 + y^2 - z^2 - 1$ (gradient $(2x, 2y,
-2z) \neq 0$ sur l'hyperboloïde, où $x^2 + y^2 = 1 + z^2 >
0$)~: dimension $2$~; la fonction de tore $G = (\rho - R)^2
+ z^2 - r^2$, $\rho = \sqrt{x^2+y^2}$, est $\mathcal
C^\infty$ près du tore ($\rho \geq R - r > 0$ là) avec
$\nabla G \neq 0$ (sa composante en $z$ est $2z$, et où $z =
0$ la composante radiale est $2(\rho - R)\ne0$ car
$\abs{\rho - R} = r$)~: dimension $2$.

(b) Pour $\varepsilon$ petit,
$(C\setminus\{0\})\cap B(0,\varepsilon)$ a exactement $2$
\omterm{def:b3:topology:components}{composantes connexes} (nappes supérieure et inférieure
épointées, chacune \omterm{def:b3:topology:pathconnected}{connexe par arcs}~: relier via cercles et
rayons). Si $C$ était une $2$-sous-variété en $0$, un
redressement donnerait un \omterm{def:b3:topology:continuity}{homéomorphisme} de $C\cap\Omega$ sur
un morceau \omterm{def:b3:topology:topology}{ouvert} d'un plan envoyant $0$ sur un point $p$~;
les petits \omterm{def:b3:topology:topology}{voisinages} épointés de $p$ dans le plan ont
\emph{une} composante, et les \omterm{def:b3:topology:continuity}{homéomorphismes} préservent le
nombre de composantes des \omterm{def:b3:topology:topology}{voisinages} épointés~:
contradiction.
\end{solution}

\begin{solution}{exo:b3:submanifolds:2}
(a) $T_aS^{n-1} = \ker\bigl(2a^{\mathsf T}\bigr) =
a^\perp$.
(b) En $p = ((R+r)\cos\theta, (R+r)\sin\theta, 0)$~: $\nabla
G = (2r\cos\theta, 2r\sin\theta, 0)$, donc le plan tangent
est $\operatorname{Vect}\bigl((-\sin\theta, \cos\theta, 0),\
(0, 0, 1)\bigr)$~: le plan vertical tangent à l'équateur
extérieur.
(c) L'hélice est une courbe plongée avec $\varphi'(t_0) =
(-\sin t_0, \cos t_0, 1) \neq 0$~: la droite tangente en
$\varphi(t_0)$ est $\varphi(t_0) + \R\,\varphi'(t_0)$, comme
le prescrit le~\cref{prop:b3:submanifolds:tangent}(2).
\end{solution}

\begin{solution}{exo:b3:submanifolds:3}
(a) $Df(x,y) = \bigl(\begin{smallmatrix}2x & -2y\\ 2y &
2x\end{smallmatrix}\bigr)$, $\det = 4(x^2 + y^2)$~: le
théorème s'applique en tout point sauf l'origine.
(b) $f(-z) = f(z)$~: jamais injective sur un ensemble
symétrique par rapport à $0$~; sur $\R^2\setminus\{0\}$ c'est
un difféomorphisme local partout mais $2$-vers-$1$
globalement. Près de $(1, 0) = f(\pm(1, 0))$, les deux
inverses sont les deux branches de racine carrée~: en
notation complexe $w \mapsto \pm\sqrt w$ (branche
principale), c'est-à-dire
\[
(u, v) \longmapsto \pm\Bigl(\sqrt{\tfrac{u +
\sqrt{u^2+v^2}}{2}},\ \frac{v}{2\sqrt{(u +
\sqrt{u^2+v^2})/2}}\Bigr).
\]
(c) Jacobien $r > 0$~: difféomorphisme local sur
$\intoo0\infty\times\R$, mais $\theta \mapsto \theta + 2\pi$
donne le même point~: localement inversible (angle déterminé
à $2\pi$ près sur un demi-plan), jamais globalement.
\end{solution}

\begin{solution}{exo:b3:submanifolds:4}
(a) $\nabla F = 3(x^2 - y,\ y^2 - x)$ s'annule ssi $y = x^2$
et $x = y^2$, c'est-à-dire $x^4 = x$~: en $(0,0)$ et $(1,1)$~;
seul $(0,0)$ est sur $\mathcal C$ ($F(1,1) = -1$). Donc sur
$\mathcal C\setminus\{0\}$, $F$ est une submersion~: une
$1$-sous-variété, localement un graphe $y = \psi(x)$ où $F_y
= 3(y^2 - x) \neq 0$ et $x = \chi(y)$ où $F_x = 3(x^2 - y)
\neq 0$ (au moins l'un des deux hors de l'origine).

(b) En $(\frac32, \frac32)$~: $\nabla F = 3(\frac94 -
\frac32)(1, 1) = \frac94(1,1)$~: droite tangente $x + y =
3$.

(c) La paramétrisation rationnelle $x = \frac{3t}{1 + t^3}$,
$y = \frac{3t^2}{1+t^3}$ passe par $0$ en $t = 0$ avec
vitesse $(3, 0)$~; l'échange $x \leftrightarrow y$ (la courbe
est symétrique, ou reparamétrer par $1/t$) donne une seconde
courbe $\mathcal C^1$ passant par $0$ de vitesse $(0, 3)$.
Deux directions tangentes indépendantes sont impossibles pour
une $1$-sous-variété (son \omterm{def:b3:submanifolds:tangent}{espace tangent} est une droite,
\cref{prop:b3:submanifolds:tangent})~: $\mathcal C$ n'est
pas une sous-variété à l'origine --- un auto-croisement
transverse.
\end{solution}

\begin{solution}{exo:b3:submanifolds:5}
(a) $F(M + H) = M^{\mathsf T}M + M^{\mathsf T}H + H^{\mathsf
T}M + H^{\mathsf T}H$~: $DF(M)(H) = M^{\mathsf T}H +
H^{\mathsf T}M$. Pour $M \in O_n$ et $S$ symétrique, $H =
\frac12MS$ donne $DF(M)(H) = \frac12(S + S^{\mathsf T}) =
S$~: surjective sur $S_n$.

(b) $O_n = F^{-1}(I)$ avec $F$ une submersion (sur $S_n$, de
dimension $\frac{n(n+1)}2$) en chacun de ses points~: une
sous-variété de dimension $n^2 - \frac{n(n+1)}2 =
\frac{n(n-1)}2$. Compacte~: fermée ($F$ \omterm{def:b3:topology:continuity}{continue}) et bornée
(les colonnes sont des vecteurs unitaires). Tangent en $I$~:
$\ker DF(I) = \{H : H + H^{\mathsf T} = 0\}$.

(c) $\bigl(\eu^{tH}\bigr)^{\mathsf T}\eu^{tH} =
\eu^{tH^{\mathsf T}}\eu^{tH} = \eu^{-tH}\eu^{tH} = I$ (la
transposition passe à travers la série~; les \omterm{thm:b3:ode:matrixexp}{exponentielles
de matrices} qui commutent se multiplient,
\cref{thm:b3:ode:matrixexp}).
\end{solution}

\begin{solution}{exo:b3:submanifolds:6}
(a) $\det(M + H) = \det M\,\det(I + M^{-1}H) = \det M\bigl(1
+ \operatorname{tr}(M^{-1}H) + O(\norm H^2)\bigr)$ pour $M$
inversible (développement de $\det$ près de $I$~: le terme
linéaire de $\prod(1 + \lambda_i)$)~; avec $\det M\cdot
M^{-1} = \operatorname{com}(M)^{\mathsf T}$~: $D\det(M)(H) =
\operatorname{tr}\bigl(\operatorname{com}(M)^{\mathsf
T}H\bigr)$, et la formule s'étend à tout $M$ par \omterm{ex:b3:lebesgue:gamma}{densité} et
continuité. Sur $SL_n$, $\det M = 1$~: $D\det(M) \ne 0$ (sa
valeur sur $H = M$ est $\operatorname{tr}(I)\cdot\dots =
n\det M = n$).

(b) $SL_n = \det^{-1}(1)$ avec $\det$ une submersion là
(valeurs dans $\R$)~: dimension $n^2 - 1$~; $T_ISL_n = \ker
D\det(I) = \{H : \operatorname{tr}H = 0\}$.

(c) $GL_n$ est un sous-ensemble \emph{\omterm{def:b3:topology:topology}{ouvert}} de $M_n(\R)$
(\cref{exo:b3:topology:8})~: une sous-variété de dimension
pleine $n^2$ (redressement~: la carte identité).
\end{solution}

\begin{solution}{exo:b3:submanifolds:7}
(a) $(y, x) = \lambda(2x, 2y)$ et $x^2 + y^2 = 1$~: $y =
2\lambda x$, $x = 2\lambda y$ donnent $x^2 = y^2 =
\frac12$. Valeurs de $xy$~: $\pm\frac12$~: maximum $\frac12$
en $\pm\frac1{\sqrt2}(1,1)$, minimum $-\frac12$ en
$\pm\frac1{\sqrt2}(1,-1)$ (l'ensemble de contrainte est
compact~: les extrema existent).

(b) Sur l'intérieur du simplexe ($p_i > 0$), Lagrange pour
$H(p) = -\sum p_i\ln p_i$ avec contrainte $\sum p_i = 1$~:
$-\ln p_i - 1 = \lambda$ pour tout $i$~: tous les $p_i$
égaux, $p_i = \frac1n$, avec $H = \ln n$. Le maximum sur le
simplexe compact est atteint~; s'il l'était sur le bord
(quelque $p_i = 0$), la distribution vit sur $\leq n - 1$
points et par récurrence $H \leq \ln(n-1) < \ln n$~: le
point critique intérieur est le maximum global ---
l'ignorance uniforme maximise l'entropie.

(c) Minimiser $(x-1)^2 + y^2$ sur l'ellipse compacte~:
$(2(x{-}1), 2y) = \lambda(\frac x2, 2y)$. Si $y \neq 0$~:
$\lambda = 1$, alors $2(x - 1) = \frac x2$ donne $x =
\frac43$, $y^2 = 1 - \frac49 = \frac59$~: distance$^2$ $=
\frac19 + \frac59 = \frac23$. Si $y = 0$~: $x = \pm2$,
distances $1$ et $3$. Points les plus proches~:
$\bigl(\frac43, \pm\frac{\sqrt5}3\bigr)$, à distance
$\sqrt{2/3} < 1$. L'équation des multiplicateurs dit que le
segment de $(1,0)$ au point le plus proche est parallèle à
$\nabla$(ellipse)~: il rencontre l'ellipse orthogonalement,
comme la géométrie l'exige.
\end{solution}

\begin{solution}{exo:b3:submanifolds:8}
Récurrence sur $n$~; $n = 1$ trivial. La fonction de Rayleigh
$f(x) = \langle Ax, x\rangle$ atteint son maximum $\lambda_1$
sur le compact $S^{n-1}$ en un certain $v_1$~; Lagrange
(\cref{thm:b3:submanifolds:lagrange}, sphère comme ensemble
de niveau) donne $2Av_1 = 2\lambda v_1$, et $\lambda =
\langle Av_1, v_1\rangle = \lambda_1$. L'hyperplan $v_1^\perp$
est $A$-invariant (symétrie~: $\langle Av, v_1\rangle =
\langle v, Av_1\rangle = 0$)~; la restriction est symétrique,
et la récurrence fournit une \omterm{def:b3:topology:basis}{base} propre orthonormée $v_2,
\dots, v_n$ de $v_1^\perp$ avec valeurs propres $\lambda_2
\geq \dots \geq \lambda_n$, chacune le maximum de $f$ sur la
sphère de l'orthocomplément restant. Courant--Fischer suit
exactement comme dans l'\cref{exo:b3:spectral:8}~: développer
$x = \sum c_iv_i$~; sur un espace-test de dimension $k$
intersecter avec $\operatorname{Vect}(v_k, \dots, v_n)$
(comptage de dimensions dans $\R^n$) pour obtenir $\min \leq
\lambda_k$, et $\operatorname{Vect}(v_1, \dots, v_k)$
réalise $\min = \lambda_k$.
\end{solution}

\begin{solution}{exo:b3:submanifolds:9}
Homothétier chaque colonne en norme $1$ divise $\abs{\det}$
par $\prod\norm{c_i}$~: il suffit de prouver $\abs{\det M}
\leq 1$ quand toutes les colonnes sont unitaires, avec
égalité ssi $M \in O_n$. La fonction $\det$ est \omterm{def:b3:topology:continuity}{continue} sur
le compact $(S^{n-1})^n$~: elle atteint un maximum $m \geq
\det I = 1 > 0$ en un certain $M$. En fixant toutes les
colonnes sauf la $i$-ème, $\det$ est linéaire en $c_i$ de
gradient la $i$-ème colonne de $\operatorname{com}(M)$~;
Lagrange sur la $i$-ème sphère~:
$\operatorname{com}(M)_{\cdot i} = \lambda_i c_i$.
L'identité $M^{\mathsf T}\operatorname{com}(M) =
\det(M)\,I$ se lit $\langle c_j,
\operatorname{com}(M)_{\cdot i}\rangle =
\det(M)\,\delta_{ij}$, c'est-à-dire $\lambda_i\langle c_j,
c_i\rangle = \det(M)\delta_{ij}$~; en prenant $j = i$~:
$\lambda_i = \det M = m \neq 0$, puis $j \neq i$ donne
$\langle c_i, c_j\rangle = 0$~: les colonnes sont
orthonormales, $M \in O_n$, $m = \abs{\det M} = 1$. Donc
$\abs{\det} \leq \prod\norm{c_i}$ toujours, avec égalité
exactement pour des colonnes orthogonales (remettre à
l'échelle)~: le volume d'un parallélépipède est \omterm{def:b3:rings:primemaximal}{maximal}, pour
des longueurs d'arêtes données, quand les arêtes sont
perpendiculaires.
\end{solution}

\begin{solution}{exo:b3:submanifolds:10}
$y = \pm\sqrt{1 - x^2}$ fonctionne près de tout point avec
$y \neq 0$~; $x = \pm\sqrt{1 - y^2}$ près de tout point avec
$x \neq 0$~; en $(1, 0)$~: le graphe $x = \sqrt{1 - y^2}$
sur $y \in \intoo{-1}1$, qui est
le~\cref{thm:b3:submanifolds:characterizations}(3) avec les
coordonnées échangées. Une certaine permutation fonctionne
toujours car la droite tangente, étant unidimensionnelle, ne
peut être simultanément verticale et horizontale --- mais
elle peut être l'une ou l'autre, donc aucun choix fixe de
coordonnée «~dépendante~» ne sert en tout point.
\end{solution}

\begin{solution}{exo:b3:submanifolds:11}
(a) L'application définissante $F(M) = M^{\mathsf T}M - I$
est \omterm{def:b3:topology:continuity}{continue}~: $O_n = F^{-1}(0)$ est fermé~; les colonnes
d'une matrice orthogonale sont des vecteurs unitaires, donc
la norme euclidienne (\omterm{thm:b3:galois:finitefields}{Frobenius}) vaut exactement $\sqrt n$~:
borné. Compact par Heine--Borel dans $M_n(\R) \cong
\R^{n^2}$.

(b) $O_n$ est l'ensemble de niveau $F = 0$ étudié dans le
chapitre~: $DF(A)H = A^{\mathsf T}H + H^{\mathsf T}A$,
surjective sur les matrices symétriques en chaque $A \in
O_n$ (étant donné $S$ symétrique, prendre $H =
\frac12AS$), donc $O_n$ est une sous-variété de dimension
$n^2 - \frac{n(n+1)}2 = \frac{n(n-1)}2$ avec
\[
T_AO_n = \ker DF(A) = \{H : A^{\mathsf T}H \text{
antisymétrique}\} = \{AK : K^{\mathsf T} = -K\} ;
\]
en $A = I$ ce sont les matrices antisymétriques.

(c) $\exp(tK)^{\mathsf T}\exp(tK) = \exp(tK^{\mathsf T})
\exp(tK) = \exp(-tK)\exp(tK) = I$ (les deux matrices $\pm
tK$ commutent, donc le produit des exponentielles est
l'exponentielle de la somme)~: $\exp(tK) \in O_n$, et
$\gamma(t) = A\exp(tK)$ est une courbe dans $O_n$ avec
$\gamma(0) = A$, $\gamma'(0) = AK$. Comme $K$ parcourt les
matrices antisymétriques, $AK$ balaye $T_AO_n$~:
l'exponentielle réalise tout l'\omterm{def:b3:submanifolds:tangent}{espace tangent} par des
courbes explicites --- le raccourci de groupe de Lie que
le~\cref{pb:b3:submanifolds:1} exploite pour $SO(3)$.
\end{solution}

\begin{solution}{exo:b3:submanifolds:12}
Existence~: intersecter $M$ avec une grande boule fermée
autour de $p$ pour obtenir un compact non vide~; la distance
\omterm{def:b3:topology:continuity}{continue} y atteint son minimum, et les points hors de la
boule sont plus loin. Condition du premier ordre~: pour une
courbe $\gamma$ dans $M$ avec $\gamma(0) = x_0$, la
fonction $h(t) = \norm{\gamma(t) - p}^2$ est différentiable
avec un minimum en $0$~:
\[
0 = h'(0) = 2\,\langle\gamma'(0),\ x_0 - p\rangle,
\]
et $\gamma'(0)$ balaye $T_{x_0}M$~: $p - x_0 \perp
T_{x_0}M$. Sphère $S(c, r)$~: l'\omterm{def:b3:submanifolds:tangent}{espace tangent} en $x_0$ est
$(x_0 - c)^\perp$, donc $p - x_0 \parallel x_0 - c$~:
$x_0$ se trouve sur la droite passant par $c$ et $p$, à
distance $r$ de $c$ --- le point sur le rayon, comme la
géométrie l'exige. Parabole~: en $x_0 = (x, x^2)$ la
tangente est engendrée par $(1, 2x)$~; l'orthogonalité à $p
- x_0 = (2 - x, -x^2)$ se lit
\[
(2 - x) - 2x^3 = 0,
\qquad\text{c'est-à-dire}\qquad 2x^3 + x - 2 = 0,
\]
avec une unique racine réelle ($x \mapsto 2x^3 + x$ est
strictement croissante) $x \approx 0.835$~; alors $x_0
\approx (0.835, 0.698)$ et $d(p, M) = \sqrt{(2 - 0.835)^2 +
0.698^2} \approx 1.358$.
\end{solution}
\begin{solution}{pb:b3:submanifolds:1}
\textbf{1.} L'application $q = t + x\mathrm i + y\mathrm j +
z\mathrm k \mapsto \bigl(\begin{smallmatrix}\alpha & \beta\\
-\bar\beta & \bar\alpha\end{smallmatrix}\bigr)$ avec $\alpha
= t + \iu x$, $\beta = y + \iu z$~: on vérifie que $1,
\mathrm i, \mathrm j, \mathrm k$ vont vers $I$,
$\bigl(\begin{smallmatrix}\iu & 0\\ 0 &
-\iu\end{smallmatrix}\bigr)$,
$\bigl(\begin{smallmatrix}0 & 1\\ -1 &
0\end{smallmatrix}\bigr)$, $\bigl(\begin{smallmatrix}0 &
\iu\\ \iu & 0\end{smallmatrix}\bigr)$, dont les produits
reproduisent la table des quaternions~: l'application est un
morphisme d'algèbres injectif, donc $\mathbb H$ hérite de
l'associativité~; $N(q) = \abs\alpha^2 + \abs\beta^2 =
\det$ est multiplicative, et la conjugaison correspond à
l'adjointe-transposée, donnant $\overline{pq} = \bar q\bar
p$. \omterm{ex:b3:groups:actions}{Centre}~: commuter avec $\mathrm i$ force $y = z = 0$,
avec $\mathrm j$ force $x = 0$~: $\R$.

\textbf{2.} $q\bar q = N(q)$~: pour $q \neq 0$, $q^{-1} =
\bar q/N(q)$~: une algèbre à division. Sur $S^3$~: $N(pq) =
1$ et $N(q^{-1}) = 1$~: un groupe~; et $S^3 \subseteq \R^4$
est la sphère unité~: une $3$-sous-variété compacte.

\textbf{3.} $v$ est pur ssi $\bar v = -v$~; alors
$\overline{qv\bar q} = q\bar v\bar q = -qv\bar q$~:
$\rho_q$ préserve $P$. La linéarité est claire~; $N(qv\bar
q) = N(q)N(v)N(q) = N(v)$~: une isométrie de $(P, N) \cong
(\R^3, \norm\cdot^2)$~: $\rho_q \in O(3)$. Et $\rho_{pq}(v)
= pqv\overline{pq} = p(qv\bar q)\bar p =
\rho_p(\rho_q(v))$~: un morphisme.

\textbf{4.} $\rho_q = \mathrm{id}$ ssi $qv = vq$ pour tout
pur $v$, ssi $q$ commute avec $\mathrm i, \mathrm j,
\mathrm k$, ssi $q$ est central (question~1)~: $q \in
\R\cap S^3 = \{\pm1\}$.

\textbf{5.} Écrire $q = t + p$ ($t \in \R$, $p$ pur)~: $1 =
N(q) = t^2 + N(p)$, donc $t = \cos\frac\theta2$ et $p =
\sin\frac\theta2\,u$ avec $N(u) = 1$ pour un certain
$\theta$ (si $p = 0$, $q = \pm1$ agit trivialement). Comme
$u^2 = -N(u) = -1$, $q$ et $u$ commutent, et $\rho_q(u) =
qu\bar q = uq\bar q = u$~: l'axe. Pour des unités pures $w
\perp u$~: la règle de produit $vw = -\langle v, w\rangle +
v\wedge w$ (développer en coordonnées depuis la table)
donne $uw = u\wedge w = -wu$. Alors
\[
\rho_q(w) = \bigl(\cos\tfrac\theta2 +
\sin\tfrac\theta2u\bigr)\,w\,\bigl(\cos\tfrac\theta2 -
\sin\tfrac\theta2u\bigr)
= \cos\theta\,w + \sin\theta\,(u\wedge w),
\]
en utilisant $uwu = -u^2w = w$ et les formules de
double-angle~: la rotation d'angle $\theta$ dans le plan
orienté $(w, u\wedge w)$.

\textbf{6.} Chaque $\rho_q$ est une rotation autour de $u$
d'angle $\theta$~: dans la \omterm{def:b3:topology:basis}{base} orthonormée $(u, w,
u\wedge w)$ sa matrice a pour déterminant $+1$~:
$\operatorname{im}\rho \subseteq SO(3)$. Surjectivité~: une
matrice $R \in SO(3)$ a $1$ comme valeur propre, car
\[
\det(R - I) = \det R\,\det(I - R^{\mathsf T}) = \det(I - R)
= (-1)^3\det(R - I),
\]
donc $\det(R - I) = 0$. Prendre un vecteur propre unitaire
$u$~; $R$ préserve $u^\perp$ et s'y restreint en une
rotation d'un certain angle $\theta$ (orthogonal plan,
déterminant $1$)~: $R = \rho_q$ pour $q =
\cos\frac\theta2 + \sin\frac\theta2\,u$. Avec la question~4
et le premier théorème d'isomorphisme
(\cref{thm:b3:groups:firstiso})~: $SO(3) \cong
S^3/\{\pm1\}$.

\textbf{7.} $O_3$ est une sous-variété compacte de
dimension~$3$ (\cref{exo:b3:submanifolds:5})~; $\det$ est
\omterm{def:b3:topology:continuity}{continue} dessus à valeurs dans $\{\pm1\}$, donc $SO(3) =
O_3\cap\{\det = 1\}$ est \omterm{def:b3:topology:topology}{ouvert} et fermé dans $O_3$~: une
union de \omterm{def:b3:topology:components}{composantes connexes}, donc elle-même une
$3$-sous-variété compacte, avec le même \omterm{def:b3:submanifolds:tangent}{espace tangent} en
$I$~: les matrices antisymétriques.

\textbf{8.} $A_u^2v = u\wedge(u\wedge v) = \langle u,
v\rangle u - v$ ($u$ unitaire), donc $A_u^3v =
u\wedge(\langle u,v\rangle u - v) = -u\wedge v$~: $A_u^3 =
-A_u$. En scindant la série exponentielle par les \omterm{def:b3:residues:singularities}{résidus} des
puissances modulo la relation $A^3 = -A$~:
\[
\eu^{\theta A_u} = I + \Bigl(\theta - \frac{\theta^3}{3!} +
\cdots\Bigr)A_u + \Bigl(\frac{\theta^2}{2!} -
\frac{\theta^4}{4!} + \cdots\Bigr)A_u^2
= I + \sin\theta\,A_u + (1 - \cos\theta)\,A_u^2 .
\]
Sur $u$~: $A_uu = 0$~: fixe. Sur $w \perp u$~:
$\eu^{\theta A_u}w = w + \sin\theta\,u\wedge w + (1 -
\cos\theta)(-w) = \cos\theta\,w + \sin\theta\,u\wedge w$~:
la rotation d'axe $u$, angle $\theta$ --- Rodrigues. Toute
rotation a cette forme (question~6)~: $\exp$ est surjective
de l'espace des matrices antisymétriques sur $SO(3)$.

\textbf{9.} Avec $q(t) = \cos\frac t2 + \sin\frac t2\,u$~:
la question~5 montre que $\rho_{q(t)}$ est la rotation
d'axe $u$ et d'angle $t$, c'est-à-dire $\rho_{q(t)} =
\eu^{tA_u}$, dont la dérivée en $t = 0$ est $A_u$. Le
quaternion court à demi-angle --- la trace analytique du
double revêtement.

\textbf{10.} $\rho_{q(t)}$ est la rotation autour de
$\mathrm k$ d'angle $t$~: en $t = 2\pi$ elle revient à
l'identité --- une boucle dans $SO(3)$. Son relèvement
vérifie $q(2\pi) = \cos\pi = -1 = -q(0)$~: le chemin relevé
n'est \emph{pas} fermé~; seulement en $t = 4\pi$ $q$
revient à $1$. Interprétation~: la boucle des rotations
complètes n'est pas contractile dans $SO(3)$ --- son
relèvement se termine sur l'autre feuille du revêtement ---
tandis que la double boucle l'est~; un corps attaché à son
environnement par des sangles (le tour de ceinture)
revient à un état non tordu après $4\pi$ mais pas après
$2\pi$. Les groupes de rotations se souviennent de la
parité des tours \omterm{def:b3:complete:complete}{complets}~; $S^3$, \omterm{def:b3:conformal:simplyconnected}{simplement connexe}, est
là où cette mémoire vit.

\textbf{11.} Quarts de tour~: $q_{\mathrm i} =
\cos\frac\pi4 + \sin\frac\pi4\,\mathrm i$, $q_{\mathrm j} =
\cos\frac\pi4 + \sin\frac\pi4\,\mathrm j$. Produit (en
appliquant d'abord le tour en $\mathrm j$)~:
\[
q_{\mathrm i}q_{\mathrm j} = \tfrac12(1 + \mathrm i)(1 +
\mathrm j) = \tfrac12\bigl(1 + \mathrm i + \mathrm j +
\mathrm k\bigr),
\]
de norme $1$, avec $\cos\frac\theta2 = \frac12$~: $\theta =
\frac{2\pi}3$, et axe $u = \frac{\mathrm i + \mathrm j +
\mathrm k}{\sqrt3}$ (partie pure normalisée). Deux quarts
de tour successifs autour d'axes orthogonaux font une
rotation de $120^\circ$ autour de la grande diagonale du
cube --- le coût de quatre multiplications réelles, l'orthogonalité
préservée exactement~: pourquoi les logiciels de vol et les
moteurs graphiques composent les rotations via les
quaternions.

\textbf{12.} Depuis la table, $\mathrm{ji} = -\mathrm k$ et
$\mathrm{ki} = \mathrm j$, donc
\[
q\,\mathrm i = a\mathrm i - b + c(\mathrm{ji}) +
d(\mathrm{ki}) = -b + a\mathrm i + d\mathrm j - c\mathrm k .
\]
En multipliant par $\bar q = a - b\mathrm i - c\mathrm j -
d\mathrm k$ avec la règle scalaire--vecteur $(t_1 +
p_1)(t_2 + p_2) = t_1t_2 - \langle p_1, p_2\rangle +
t_1p_2 + t_2p_1 + p_1\wedge p_2$, où $p_1 = (a, d, -c)$ et
$p_2 = (-b, -c, -d)$~: la partie scalaire est $-ab - (-ab)
= 0$ (pur, comme il se doit), et la partie vectorielle est
\[
b(b, c, d) + a(a, d, -c) + (-c^2 - d^2,\ bc + ad,\ bd - ac)
= \bigl(a^2 + b^2 - c^2 - d^2,\ 2(bc + ad),\ 2(bd -
ac)\bigr) :
\]
la première colonne de $R_q$. L'application cyclique
$\sigma(\mathrm i) = \mathrm j$, $\sigma(\mathrm j) =
\mathrm k$, $\sigma(\mathrm k) = \mathrm i$ préserve les
relations $\mathrm i^2 = \mathrm j^2 = \mathrm k^2 =
\mathrm{ijk} = -1$ (le mot $\mathrm{ijk}$ est
cycliquement invariant à la relation $\mathrm{ijk} =
\mathrm{jki}$ près, qui vaut dans tout anneau~: conjuguer
$\mathrm{ijk} = -1$ par l'inversible $\mathrm i$), donc
$\sigma$ s'étend en un automorphisme de $\R$-algèbre, et
$\sigma(\rho_q(v)) = \rho_{\sigma(q)}(\sigma(v))$. En
déroulant, l'image de $\mathrm j$ est la formule de la
première colonne après la substitution $(b, c, d) \to (c,
d, b)$ avec la \omterm{def:b3:topology:basis}{base} renommée $\mathrm i \to \mathrm j \to
\mathrm k \to \mathrm i$, qui est exactement la seconde
colonne affichée~; un tour de plus donne la troisième.

\textbf{13.} En sommant la diagonale, $\operatorname{tr}R_q
= 3a^2 - (b^2 + c^2 + d^2) = 4a^2 - 1$ (norme unitaire), et
avec $a = \cos\frac\theta2$~: $4\cos^2\frac\theta2 - 1 = 1 +
2\cos\theta$. Partie antisymétrique~: les trois entrées
indépendantes de $R_q - R_q^{\mathsf T}$ sont $4ab, 4ac,
4ad$ (aux positions $(3,2), (1,3), (2,1)$), donc
$\frac12(R_q - R_q^{\mathsf T}) = A_m$ avec $m =
2a\,(b, c, d) = 2\cos\frac\theta2\sin\frac\theta2\,u =
\sin\theta\,u$. Algorithme~: $\theta =
\arccos\frac{\operatorname{tr}R - 1}{2} \in \intcc0\pi$~;
si $0 < \theta < \pi$, lire $u$ sur $\frac{R - R^{\mathsf
T}}{2\sin\theta}$ et poser $q = \pm(\cos\frac\theta2 +
\sin\frac\theta2 u)$~; si $\theta = 0$, $q = \pm1$. Pour
$\theta = \pi$~: $a = 0$, et Rodrigues (question~8) donne
$R = I + 2A_u^2 = I + 2(uu^{\mathsf T} - I) =
2uu^{\mathsf T} - I$, c'est-à-dire $R + I = 2uu^{\mathsf
T}$~; toute colonne non nulle de $R + I$, normalisée, est
$\pm u$, et $q = \pm u$.

\textbf{14.} Avec $a = b = c = d = \frac12$~: toutes les
entrées diagonales s'annulent, $2(bc - ad) = 0$, $2(bd +
ac) = 1$, $2(bc + ad) = 1$, $2(cd - ab) = 0$, $2(bd - ac) =
0$, $2(cd + ab) = 1$~:
\[
R_q = \begin{pmatrix} 0 & 0 & 1\\ 1 & 0 & 0\\ 0 & 1 & 0
\end{pmatrix},
\]
la permutation cyclique $e_1 \to e_2 \to e_3 \to e_1$. Sa
trace est $0 = 1 + 2\cos\theta$, donc $\theta =
\frac{2\pi}3$, et elle fixe $(1,1,1)$~: la rotation de
$120^\circ$ autour de la grande diagonale --- précisément
la réponse de la question~11, maintenant visible comme la
matrice qui cycle les axes de coordonnées.

\textbf{15.} Sur la \omterm{def:b3:topology:basis}{base} on vérifie $\Phi(\bar q) =
\Phi(q)^*$ (la matrice de $\bar q$ a $\alpha' =
\bar\alpha$, $\beta' = -\beta$, qui est la conjuguée
transposée de $\bigl(\begin{smallmatrix}\alpha & \beta\\
-\bar\beta & \bar\alpha\end{smallmatrix}\bigr)$). D'où
$\Phi(q)^*\Phi(q) = \Phi(\bar qq) = N(q)I$ et $\det\Phi(q)
= \abs\alpha^2 + \abs\beta^2 = N(q)$~: pour $q \in S^3$,
$\Phi(q) \in SU(2)$, et $\Phi$ est un morphisme injectif
(question~1). Surjectivité~: soit $U =
\bigl(\begin{smallmatrix}\alpha & \beta\\ \gamma &
\delta\end{smallmatrix}\bigr)$ avec $\det U = 1$~; alors
$U^{-1} = \bigl(\begin{smallmatrix}\delta & -\beta\\
-\gamma & \alpha\end{smallmatrix}\bigr)$, et $U^{-1} = U^*
= \bigl(\begin{smallmatrix}\bar\alpha & \bar\gamma\\
\bar\beta & \bar\delta\end{smallmatrix}\bigr)$ force
$\delta = \bar\alpha$, $\gamma = -\bar\beta$, puis $1 =
\det U = \abs\alpha^2 + \abs\beta^2$~: $U = \Phi(q)$ pour
le quaternion unitaire $q$ de coordonnées $\alpha = a +
\iu b$, $\beta = c + \iu d$. Donc $S^3 \cong SU(2)$.

\textbf{16.} Les scalaires réels sont centraux et $N(p) =
1$ donne $\overline{pq\bar p} = p\bar q\bar p$, donc
$pq\bar p + \overline{pq\bar p} = p(q + \bar q)\bar p = q
+ \bar q$~: la partie réelle est invariante.
Réciproquement soient $\operatorname{Re}q =
\operatorname{Re}q' = a$~; alors les parties pures ont la
même norme $\sqrt{1 - a^2} = s$. Si $s = 0$, $q = q' =
\pm1$. Si $s > 0$, écrire $q = a + su$, $q' = a + su'$ avec
$u, u'$ purs unitaires~; la question~6 fournit une
rotation portant $u'$ sur $u$, c'est-à-dire $p \in S^3$
avec $\rho_p(u') = u$, puis $pq'\bar p = a +
s\rho_p(u') = q$. Les classes de $S^3$ sont donc $\{1\}$,
$\{-1\}$, et pour chaque $a \in \intoo{-1}1$ la $2$-sphère
$\{a + su : u \text{ pur unitaire}\}$ de rayon $s$. Sous
$\Phi$, $\operatorname{tr}\Phi(q) = \alpha + \bar\alpha =
2\operatorname{Re}q$~: les classes de $SU(2)$ sont les
ensembles de niveau de la trace. En projetant~: si $q' =
pq\bar p$ alors $\rho_{q'} = \rho_p\rho_q\rho_p^{-1}$~;
réciproquement $\rho_{q'} = \rho_p\rho_q\rho_p^{-1} =
\rho_{pq\bar p}$ force $q' = \pm pq\bar p$ (noyau), donc
$\operatorname{Re}q' = \pm\operatorname{Re}q$,
c'est-à-dire $\cos\frac{\theta'}2 =
\abs{\operatorname{Re}q'} = \abs{\operatorname{Re}q} =
\cos\frac\theta2$ pour les angles dans $\intcc0\pi$~: des
rotations conjuguées ont des angles égaux.
Réciproquement, des angles égaux permettent des
représentants de même partie réelle non négative, conjugués
par ce qui précède~: dans $SO(3)$, la \omterm{ex:b3:groups:actions}{classe de conjugaison}
d'une rotation est exactement son angle.

\textbf{17.} $N$ est multiplicative, donc $\abs\cdot$ est
une norme multiplicative sur $\mathbb H \cong \R^4$ et
$\abs{q^k} = \abs q^k$~: la série $\sum q^k/k!$ converge
absolument dans l'espace de dimension finie (donc
\omterm{def:b3:complete:complete}{complet}), dominée par $\sum\abs q^k/k! = \eu^{\abs q}$.
Pour un pur unitaire $u$~: $u^2 = -1$, donc $(\theta
u)^{2m} = (-1)^m\theta^{2m}$ et $(\theta u)^{2m+1} =
(-1)^m\theta^{2m+1}u$~; en scindant la série,
\[
\exp(\theta u) = \sum_m\frac{(-1)^m\theta^{2m}}{(2m)!} +
u\sum_m\frac{(-1)^m\theta^{2m+1}}{(2m+1)!} = \cos\theta +
\sin\theta\,u .
\]
Tout $q \in S^3$ est $\cos\alpha + \sin\alpha\,u$ avec
$\alpha \in \intcc0\pi$ (question~5)~: $q = \exp(\alpha
u)$, donc $\exp(P) = S^3$. Enfin $\exp(su) = \cos s + \sin
s\,u = q(2s)$ dans la notation de la question~9, et
$\rho_{q(t)} = \eu^{tA_u}$ là~: $\rho_{\exp(su)} =
\eu^{2sA_u}$.

\textbf{18.} En développant $vw$ coordonnée par coordonnée
avec la table~: les produits $\mathrm i\cdot\mathrm i =
-1$, \dots\ donnent le scalaire $-(v_1w_1 + v_2w_2 +
v_3w_3)$, et les produits mixtes ($\mathrm{ij} = \mathrm
k$, $\mathrm{ji} = -\mathrm k$, \dots) donnent le vecteur
$(v_2w_3 - v_3w_2,\ v_3w_1 - v_1w_3,\ v_1w_2 - v_2w_1)$~:
$vw = -\langle v, w\rangle + v\wedge w$. En soustrayant le
produit inversé~: $vw - wv = 2\,v\wedge w$ (les parties
scalaires s'annulent, les produits vectoriels
s'additionnent). Pour les matrices, avec $a\wedge(b\wedge
c) = b\langle a, c\rangle - c\langle a, b\rangle$~:
\[
[A_v, A_w]x = v\wedge(w\wedge x) - w\wedge(v\wedge x)
= w\langle v, x\rangle - v\langle w, x\rangle
= (v\wedge w)\wedge x = A_{v\wedge w}x .
\]
Dérivée~: $\overline{\exp(tv)} = \exp(-tv)$ (la conjugaison
est \omterm{def:b3:topology:continuity}{continue} et nie les quaternions purs), donc
\[
\frac{\dd}{\dd t}\Bigr|_{t=0}\exp(tv)\,x\,\exp(-tv) = vx -
xv = 2\,v\wedge x = 2A_vx :
\]
la différentielle est $v \mapsto 2A_v$, une bijection
linéaire de $P$ sur les matrices antisymétriques, et $[2A_v,
2A_w] = 4A_{v\wedge w} = 2A_{2v\wedge w} =
2A_{[v,w]}$ montre qu'elle transporte le \omterm{def:b3:groups:derived}{commutateur}
quaternionique vers le \omterm{def:b3:groups:derived}{commutateur} matriciel.

\textbf{19.} En appliquant $\rho$~: $\rho(\varepsilon(t)) =
\rho(s(R(t)))\,\rho(q(t))^{-1} = R(t)R(t)^{-1} =
\operatorname{id}$, donc $\varepsilon(t) \in \ker\rho =
\{\pm1\}$ (question~4). Comme produit des \omterm{def:b3:topology:continuity}{applications
continues} $t \mapsto s(R(t))$ et $t \mapsto q(t)^{-1} =
\bar{q(t)}$, $\varepsilon$ est \omterm{def:b3:topology:continuity}{continue} sur l'intervalle
\omterm{def:b3:topology:connected}{connexe} $\intcc0{2\pi}$ à valeurs dans la paire discrète
$\{\pm1\}$~: elle est constante, disons $\varepsilon(t)
\equiv \varepsilon$. Mais $R(0) = R(2\pi) = I$, donc
$s(R(0)) = s(R(2\pi))$, tandis que $s(R(0)) =
\varepsilon\,q(0) = \varepsilon$ et $s(R(2\pi)) =
\varepsilon\,q(2\pi) = -\varepsilon$~: contradiction. Aucune
section \omterm{def:b3:topology:continuity}{continue} n'existe~: l'ambiguïté de signe $\pm q$
est globale, non un défaut d'une formule particulière.

\textbf{20.} \emph{Surjectivité~:} tout $R \in SO(3)$ est
$\eu^{\theta A_u}$ pour un certain $u$ unitaire et $\theta
\in \intcc0{2\pi}$ (questions~6 et~8)~; si $\theta > \pi$,
Rodrigues donne $\eu^{\theta A_u} = \eu{(2\pi -
\theta)A_{-u}}$ (tous deux valent $I + \sin\theta A_u + (1
- \cos\theta)A_u^2$, et $A_{-u} = -A_u$ avec $\sin(2\pi -
\theta) = -\sin\theta$, $\cos(2\pi - \theta) =
\cos\theta$), donc $R = E(v)$ avec $\norm v \leq \pi$.
\emph{Injectivité à l'intérieur~:} si $E(v) = E(v') \neq I$
avec $\norm v, \norm{v'} < \pi$, la question~13 récupère le
même angle $\theta = \norm v = \norm{v'} \in \intoo0\pi$
depuis la trace et, comme $\sin\theta \neq 0$, le même axe
depuis la partie antisymétrique~: $v = v'$~; et $E(v) = I$
force $\theta \in \{0\}$ sur la boule ouverte. \emph{Bord~:}
$E(\pi u) = I + 2A_u^2 = 2uu^{\mathsf T} - I$ ne dépend de
$u$ qu'à travers $uu^{\mathsf T}$, d'où $E(\pi u) =
E(-\pi u)$~; réciproquement $2uu^{\mathsf T} - I =
2u'u'^{\mathsf T} - I$ appliqué à $u$ donne $u = \langle
u', u\rangle u'$, donc $u' = \pm u$. Intérieur et bord ne se
rencontrent jamais (trace $> -1$ contre $= -1$). Donc $E$
induit une bijection \omterm{def:b3:topology:continuity}{continue} de la
boule-avec-collage-antipodale --- compacte --- sur $SO(3)$~:
un \omterm{def:b3:topology:continuity}{homéomorphisme}, et $SO(3) \cong \mathbb{RP}^3$. Un
diamètre de $\pi u$ à $-\pi u$ a des extrémités collées~:
c'est une boucle dans $SO(3)$, et sa description par $E$
correspond à la famille de rotations autour de $u$ de la
question~10 balayant un tour \omterm{def:b3:complete:complete}{complet}.

\textbf{21.} $\rho$ est un morphisme et $\rho_p^{-1} =
\rho_{p^{-1}} = \rho_{\bar p}$, donc
$\rho_p\rho_q\rho_p^{-1} = \rho_{pq\bar p}$~; par la
question~16, conjuguer une rotation préserve son angle et
fait tourner son axe par $\rho_p$. Involutions~: $\rho_q^2
= \rho_{q^2} = \operatorname{id}$ ssi $q^2 = \pm1$. Si $q^2
= 1$ alors $(q - 1)(q + 1) = q^2 - 1 = 0$ (scalaires
centraux, donc cette factorisation est valide) et $q =
\pm1$ dans l'anneau à division $\mathbb H$, donnant
$\rho_q = I$, exclu~; $q^2 = -1$ avec $q = a + su$ donne
$a^2 - s^2 + 2as\,u = -1$, donc $a = 0$, $s = 1$~: $q$ est
un pur unitaire $w$, et $\rho_w$ est le demi-tour autour de
$w$ (angle $\pi$, question~5). \omterm{ex:b3:groups:actions}{Centre}~: si $\rho_q$ commute
avec tout $\rho_p$, alors $\rho_{pq\bar p} = \rho_q$, donc
$pq\bar p = \varepsilon(p)\,q$ avec $\varepsilon(p) \in
\{\pm1\}$~; $p \mapsto \varepsilon(p) = (pq\bar
p)q^{-1}$ est \omterm{def:b3:topology:continuity}{continue} sur le $S^3$ \omterm{def:b3:topology:connected}{connexe} et vaut $1$ en
$p = 1$, d'où $\equiv 1$~: $q$ commute avec tout $S^3$,
donc avec tout $\mathbb H$ (remettre à l'échelle), donc $q
\in \R \cap S^3 = \{\pm1\}$ (question~1) et $\rho_q = I$~:
le \omterm{ex:b3:groups:actions}{centre} est trivial.

\textbf{22.} Comme $u \perp w$ sont des purs unitaires, $uw
= u\wedge w$ est pur (question~18), donc
\[
w' = qw = \cos\tfrac\theta2\,w + \sin\tfrac\theta2\,u\wedge
w
\]
est pur, de norme $\abs q\abs w = 1$. Alors $w'w = qw\cdot
w = qw^2 = -q$, et
\[
\rho_{w'}\rho_w = \rho_{w'w} = \rho_{-q} = \rho_q .
\]
Les deux axes $w$ et $w' = \cos\frac\theta2 w +
\sin\frac\theta2(u\wedge w)$ se trouvent dans le plan
$u^\perp$ orthogonal à l'axe de rotation, et $\langle w',
w\rangle = \cos\frac\theta2$~: ils font le demi-angle
$\frac\theta2$. C'est la génération classique~: deux
demi-tours autour d'axes se rencontrant à l'angle
$\frac\theta2$ se composent en la rotation d'angle $\theta$
autour de leur perpendiculaire commune.

\textbf{23.} La compacité est la question~7. Connexité par
arcs~: $SO(3) = \rho(S^3)$ est l'image \omterm{def:b3:topology:continuity}{continue} de la
sphère \omterm{def:b3:topology:pathconnected}{connexe par arcs}~; alternativement, pour $R =
\eu^{A}$ avec $A$ antisymétrique (question~8), $t \mapsto
\eu^{tA}$ est un chemin dans $SO(3)$ de $I$ à $R$
(orthogonal car $(\eu^{tA})^{\mathsf T} = \eu^{-tA}$,
déterminant $1$ par continuité depuis $t = 0$). Pour
$O(3)$~: $\det$ est \omterm{def:b3:topology:continuity}{continue} sur $\{\pm1\}$, donc $O(3)$
est disconnexe, $O(3) = SO(3) \sqcup D\,SO(3)$ pour tout
$D$ fixe avec $\det D = -1$ (p.ex.\ $D = -I$), et la
multiplication à gauche par $D$ est un \omterm{def:b3:topology:continuity}{homéomorphisme}~:
exactement deux composantes, chacune une copie de $SO(3)$.
Le $\rho$ deux-vers-un, sans section par la question~19,
est ainsi un double revêtement honnête d'un groupe compact
\omterm{def:b3:topology:connected}{connexe} par le $S^3$ \omterm{def:b3:conformal:simplyconnected}{simplement connexe} --- la géométrie
derrière le tour de ceinture.

\textbf{24.} (i) Matrices~:
\[
R_1 = \begin{pmatrix} 0 & -1 & 0\\ 1 & 0 & 0\\ 0 & 0 & 1
\end{pmatrix},
\quad
R_2 = \begin{pmatrix} 1 & 0 & 0\\ 0 & 0 & -1\\ 0 & 1 & 0
\end{pmatrix},
\quad
R_2R_1 = \begin{pmatrix} 0 & -1 & 0\\ 0 & 0 & -1\\ 1 & 0 &
0\end{pmatrix}.
\]
Trace $0 = 1 + 2\cos\theta$ donne $\cos\theta = -\frac12$~:
$\theta = \frac{2\pi}3$. Partie antisymétrique $\frac{R -
R^{\mathsf T}}2$ a des entrées encodant $\sin\theta\,(v_3,
-v_2, v_1)$-wise l'axe~: ici $\frac{R - R^{\mathsf T}}2 =
\frac12\bigl(\begin{smallmatrix}0 & -1 & -1\\ 1 & 0 & -1\\ 1
& 1 & 0\end{smallmatrix}\bigr)$, qui se lit (dictionnaire
de la Partie~V $A_v$) $v\sin\theta = \frac12(1, -1, 1)$~;
avec $\sin\frac{2\pi}3 = \frac{\sqrt3}2$~: $v =
\frac1{\sqrt3}(1, -1, 1)$.
(ii) Quaternions~: $q_1 = \cos\frac\pi4 + \sin\frac\pi4\,k
= \frac{\sqrt2}2(1 + k)$, $q_2 = \frac{\sqrt2}2(1 + i)$, et
\[
q_2q_1 = \tfrac12(1 + i)(1 + k) = \tfrac12(1 + k + i + ik)
= \tfrac12\bigl(1 + i - j + k\bigr)
\]
($ik = -j$). Donc $\cos\frac\theta2 = \frac12$~: $\theta =
\frac{2\pi}3$, et la partie vectorielle $\frac12(i - j +
k)$ a pour direction $\frac1{\sqrt3}(1, -1, 1)$ --- la même
réponse, la voie quaternionique ne demandant qu'une ligne
de multiplication au lieu d'un produit matriciel~: la
raison pratique pour laquelle les logiciels de vol
composent les attitudes dans $S^3$.

\textbf{25.} $I + K$ inversible~: $(I + K)v = 0$ donne $0 =
\langle v, v\rangle + \langle v, Kv\rangle = \norm v^2$
(l'antisymétrie tue le second terme)~: $v = 0$.
Orthogonalité de $C = C(K)$~: en utilisant $(I \pm
K)^{\mathsf T} = I \mp K$ et le fait que les quatre
matrices $I \pm K$, $(I \pm K)^{-1}$ commutent (expressions
polynomiales en $K$, plus limites)~:
\[
C^{\mathsf T}C = (I + K)^{-\mathsf T}(I - K)^{\mathsf T}
(I - K)(I + K)^{-1}
= (I - K)^{-1}(I + K)(I - K)(I + K)^{-1} = I .
\]
Déterminant~: $\det(I - K) = \det\bigl((I - K)^{\mathsf
T}\bigr) = \det(I + K)$, donc $\det C = 1$~: $C \in
SO(n)$. Pas de valeur propre $-1$~: $Cv = -v$ signifie $(I
- K)w = -(I + K)w$ pour $w = (I + K)^{-1}v$, c'est-à-dire
$2w = 0$~: $v = 0$. Inversion~: depuis $C(I + K) = I - K$,
résoudre $K(I + C) = I - C$~; comme $-1 \notin
\operatorname{Sp}C$, $I + C$ est inversible et $K = (I -
C)(I + C)^{-1}$, qui est antisymétrique dès que $C$ est
orthogonale sans valeur propre $-1$ (transposer
l'expression et utiliser $C^{\mathsf T} = C^{-1}$~:
$K^{\mathsf T} = (I - C^{-1})(I + C^{-1})^{-1} = (C -
I)(C + I)^{-1} = -K$). Les deux applications sont
mutuellement inverses par construction~: une
paramétrisation rationnelle globale du morceau \omterm{def:b3:topology:topology}{ouvert} \omterm{def:b3:topology:interior}{dense}
de $SO(n)$ évitant la valeur propre $-1$ --- pas de série,
pas de trigonométrie, et en dimension~$3$ c'est la
substitution de demi-angle $K = \tan\frac\theta2\,A_v$ en
déguisement.
\end{solution}
